% Section 15.1: The dictionary in $\R^3$ | Part: 2 | Depends on: sec:fields, sec:lie, sec:ga-calculus | Status: done | Original lines: 2375-2415
%!TEX root = ../../main.tex
\subsection{The dictionary in $\R^3$}

% REWRITTEN: definitions of vol and of the flux 2-form added (both are used in the table before Part III).
In the table, $\vol=\dd x\wedge\dd y\wedge\dd z$ is the \emph{volume form} of $\R^3$ (it gives volume $1$ to the unit cube), and the
\emph{flux $2$-form} of a vector field $F$ is $\iota_F\vol$, computed in \eqref{eq:iotavol}; Section~\ref{sec:flux} explains why its
integral over a surface is the flux of $F$. The Hodge star on forms is the one of Section~\ref{sec:duality}.

\begin{table}[ht]
\centering
\small
\renewcommand{\arraystretch}{1.4}
\begin{tabular}{@{}p{3.0cm}p{3.4cm}p{4.2cm}p{4.0cm}@{}}
\toprule
\textbf{Object} & \textbf{Vector calculus} & \textbf{Forms} & \textbf{Geometric algebra}\\
\midrule
gradient & $\operatorname{grad}f$ & $\dd f=(\operatorname{grad}f)^\flat$ & $\vd f$\\
curl & $\operatorname{curl}F$ & $\dd F^\flat=\iota_{\operatorname{curl}F}\vol=\hodge(\operatorname{curl}F)^\flat$ & $\vd\wedge F=\pss\,\operatorname{curl}F$\\
divergence & $\operatorname{div}F$ & $\dd(\iota_F\vol)=(\operatorname{div}F)\vol$ & $\vd\lc F=\vd\cdot F$\\
flux $2$-form & $F\cdot n\,\dd A$ & $\iota_F\vol=\hodge F^\flat$ & $F\pss=\pss F$\\
cross product & $a\times b$ & $\bigl(\hodge(a^\flat\wedge b^\flat)\bigr)^\sharp$ & $(a\wedge b)\pss^{-1}$\\
$\dd\dd=0$ & $\operatorname{curl}\operatorname{grad}=0,\ \operatorname{div}\operatorname{curl}=0$ & $\dd\dd=0$ & $\vd\wedge\vd\wedge=0$\\
all first derivatives at once & --- & $\dd$ and $\codiff$ separately & $\vd F=\vd\lc F+\vd\wedge F$\\
\bottomrule
\end{tabular}
\caption{The calculus of $\R^3$ in three languages (Euclidean metric, orientation $\pss=e_1e_2e_3$, $\pss^2=-1$; forms and multivectors identified by $\dd x^i\leftrightarrow e_i$).}
\label{tab:dictionary3}
\end{table}

\begin{proposition}[three-language identities]\label{prop:threelang}
For a vector field $F=F^ie_i$ on $\R^3$ (identified with $F^\flat=F_i\,\dd x^i$):
\begin{enumerate}[label=(\roman*)]
\item $\iota_F\vol=\hodge F^\flat\ \leftrightarrow\ F\pss$, with components $F^1\dd y\wedge\dd z+F^2\dd z\wedge\dd x+F^3\dd x\wedge\dd y\ \leftrightarrow\ F^1e_2e_3+F^2e_3e_1+F^3e_1e_2$;
\item $\dd F^\flat=\iota_{\operatorname{curl}F}\vol\ \leftrightarrow\ \vd\wedge F=\pss\operatorname{curl}F$;
\item $\dd(\iota_F\vol)=(\operatorname{div}F)\vol\ \leftrightarrow\ \vd\wedge(F\pss)=(\vd\cdot F)\,\pss$.
\end{enumerate}
\end{proposition}
\begin{proof}
(i) By Theorem~\ref{thm:hodgeGA}, $\hodge F=\rev F\pss=F\pss$ for a vector; $e_1\pss=e_2e_3$, $e_2\pss=e_3e_1$, $e_3\pss=e_1e_2$ (e.g.\ $e_2e_1e_2e_3=-e_1e_2e_2e_3=-e_1e_3=e_3e_1$). Compare with the formula \eqref{eq:iotavol} for $\iota_F\vol$ of Section~\ref{sec:interior}. (ii) $\vd\wedge F=\sum_{i<j}(\partial_iF_j-\partial_jF_i)\,e_ie_j=(\partial_2F_3-\partial_3F_2)e_2e_3+(\partial_3F_1-\partial_1F_3)e_3e_1+(\partial_1F_2-\partial_2F_1)e_1e_2$, which is $\pss\operatorname{curl}F$ by the same table of $e_i\pss$; the forms side is the computation of $\dd F^\flat$ in Section~\ref{sec:dform}.
% REWRITTEN (F5): both references now point backward (sec:interior, sec:dform) instead of forward (sec:intk, sec:stokes). (iii) Expanding, $\vd\wedge(F^1e_2e_3+F^2e_3e_1+F^3e_1e_2)=(\partial_1F^1+\partial_2F^2+\partial_3F^3)e_1e_2e_3$, since $e_1\wedge e_2e_3=e_2\wedge e_3e_1=e_3\wedge e_1e_2=\pss$.
\end{proof}

\begin{example}[all three languages on one field]\label{ex:threelang}
Let $F=(-y,\,x,\,z)$ on $\R^3$. \emph{Vector calculus:} $\operatorname{curl}F=(0,0,2)$, $\operatorname{div}F=1$.
\emph{Forms:} $F^\flat=-y\,\dd x+x\,\dd y+z\,\dd z$, so $\dd F^\flat=-\dd y\wedge\dd x+\dd x\wedge\dd y=2\,\dd x\wedge\dd y=\iota_{(0,0,2)}\vol$ \checkmark; and $\iota_F\vol=-y\,\dd y\wedge\dd z+x\,\dd z\wedge\dd x+z\,\dd x\wedge\dd y$, whose $\dd$ is $\partial_zz\,\dd z\wedge\dd x\wedge\dd y=\vol$ (the other two terms contribute $\partial_x(-y)=\partial_y(x)=0$), i.e.\ $\operatorname{div}F=1$ \checkmark.
\emph{Geometric algebra:} the single product $\vd F=\vd\lc F+\vd\wedge F=1+\pss(0,0,2)=1+2\,e_1e_2$ contains both answers at once: scalar part $\operatorname{div}F=1$, bivector part $\pss\operatorname{curl}F=2e_1e_2$.
\end{example}
